
\documentclass{article}
\usepackage{colortbl}
\usepackage{makecell}
\usepackage{multirow}
\usepackage{supertabular}

\begin{document}

\newcounter{utterance}

\centering \large Interaction Transcript for game `dond', experiment `semi\_en', episode 5 with reflexion{-}llama{-}3.1{-}8b{-}instruct.
\vspace{24pt}

{ \footnotesize  \setcounter{utterance}{1}
\setlength{\tabcolsep}{0pt}
\begin{supertabular}{c@{$\;$}|p{.15\linewidth}@{}p{.15\linewidth}p{.15\linewidth}p{.15\linewidth}p{.15\linewidth}p{.15\linewidth}}
    \# & $\;$A & \multicolumn{4}{c}{Game Master} & $\;\:$B\\
    \hline

    \theutterance \stepcounter{utterance}  
    & & \multicolumn{4}{p{0.6\linewidth}}{
        \cellcolor[rgb]{0.9,0.9,0.9}{
            \makecell[{{p{\linewidth}}}]{
                \texttt{\tiny{[P1$\langle$GM]}}
                \texttt{You have attempted this same task before and did not succeed. These are the plans you wrote after those attempts.} \\
\texttt{Trial \#1: New plan: } \\
\\ 
\texttt{In my last attempt, I was able to make a trade that seemed beneficial, but ultimately, I failed to achieve the target state of the game. Upon reflection, I realize that I didn't consider the other player's secret value function and my own value function in a more strategic way. } \\
\\ 
\texttt{In the beginning, I tried to gauge the other player's interest in the hammer, which is a valuable item for me. However, I didn't use this information effectively to make a proposal that would maximize my score. } \\
\\ 
\texttt{Next time, I will start by making a proposal that takes into account both my own value function and the other player's interest in the items. I will also try to identify the items that are most valuable to both of us and make a trade that benefits both parties. } \\
\\ 
\texttt{For example, I could propose giving the other player the hammer and one rope in exchange for the pepper and one rope. This way, I would get the hammer, which is highly valued by me, and the other player would get the items that are most valuable to them. } \\
\\ 
\texttt{By considering the other player's secret value function and making a proposal that benefits both parties, I hope to achieve a higher score and win the game.} \\
\\ 
\texttt{Trial \#2: New plan: } \\
\\ 
\texttt{In my last attempt, I was able to make a trade that seemed beneficial, but ultimately, I failed to achieve the target state of the game. Upon reflection, I realize that I didn't consider the other player's secret value function and my own value function in a more strategic way. } \\
\\ 
\texttt{In the beginning, I tried to gauge the other player's interest in the hammer, which is a valuable item for me. However, I didn't use this information effectively to make a proposal that would maximize my score. } \\
\\ 
\texttt{Moreover, I made a mistake by not making a secret proposal at the end of the game. I should have proposed a deal that would have maximized my score, rather than agreeing to the other player's proposal. } \\
\\ 
\texttt{Next time, I will start by making a proposal that takes into account both my own value function and the other player's interest in the items. I will also try to identify the items that are most valuable to both of us and make a trade that benefits both parties. } \\
\\ 
\texttt{To avoid making a secret proposal too late, I will make sure to keep track of the items that are left and the value of each item. This way, I can make a proposal that maximizes my score at the end of the game. } \\
\\ 
\texttt{For example, I could propose giving the other player the items that are most valuable to them, in exchange for the items that are most valuable to me. This way, I would get the items that I need to maximize my score, and the other player would get the items that they need to maximize their score. } \\
\\ 
\texttt{By considering the other player's secret value function and making a proposal that benefits both parties, I hope to achieve a higher score and win the game.} \\
\\ 
\\ 
\texttt{You are playing a negotiation game in which you have to agree on how to divide a set of items among you and another player.} \\
\\ 
\texttt{Rules:} \\
\texttt{(a) You and the other player are given a set of items. Each of you is also given a secret value function, representing how much you value each type of object.} \\
\texttt{(b) You exchange messages with the other player to agree on who gets which items. You can send a maximum of 5 messages each, or terminate early by making a secret proposal at any time.} \\
\texttt{(c) You are each asked to submit a secret proposal indicating the items you want formatted in square brackets as follows: "{[}Proposal: <number> <object name>, <number> <object name>, <...>{]}"} \\
\texttt{(d) If your proposals are complementary, i.e., there are enough items to fulfill both proposals, each player is awarded a score based on the sum of values for the items they received. Otherwise, both of you get zero points.} \\
\texttt{(e) Your goal is to maximize the score you receive.} \\
\\ 
\texttt{Let us start.} \\
\\ 
\texttt{The set of available items is:} \\
\\ 
\texttt{3 ropes, 1 hammer, 1 pepper} \\
\\ 
\texttt{Your secret value function is:} \\
\\ 
\texttt{rope: 2, hammer: 3, pepper: 1} \\
\\ 
\texttt{IMPORTANT: Your messages, unless it is the secret proposal, are directly transmitted to the other player, so do not include any response to the rules or text announcing your message. To make a secret proposal, use the indicated format. Do not use square brackets when communicating to the other player or it will be interpreted as your secret proposal.} \\
            }
        }
    }
    & \\ \\

    \theutterance \stepcounter{utterance}  
    & \multicolumn{4}{p{0.6\linewidth}}{
        \cellcolor[rgb]{0.9,0.9,0.9}{
            \makecell[{{p{\linewidth}}}]{
                \texttt{\tiny{[P1$\rangle$GM]}}
                \texttt{I'd like to start by making a proposal that takes into account both my own value function and the other player's interest in the items. I'll propose giving the other player the hammer and one rope in exchange for the pepper and one rope. } \\
\\ 
\texttt{hammer: 3, rope: 2, pepper: 1} \\
\\ 
\texttt{I think this proposal is a good starting point because it gives the other player the hammer, which is likely to be highly valued by them, and I get the pepper, which is the least valuable item to me. I also get to keep one rope, which is valuable to me. } \\
\\ 
\texttt{What do you think?} \\
            }
        }
    }
    & & \\ \\

    \theutterance \stepcounter{utterance}  
    & & & \multicolumn{2}{p{0.3\linewidth}}{
        \cellcolor[rgb]{0.9,0.9,0.9}{
            \makecell[{{p{\linewidth}}}]{
                \texttt{\tiny{[GM$|$GM]}}
                \texttt{continue} \\
            }
        }
    }
    & & \\ \\

    \theutterance \stepcounter{utterance}  
    & & & \multicolumn{2}{p{0.3\linewidth}}{
        \cellcolor[rgb]{0.9,0.9,0.9}{
            \makecell[{{p{\linewidth}}}]{
                \texttt{\tiny{[GM$|$GM]}}
                \texttt{I'd like to start by making a proposal that takes into account both my own value function and the other player's interest in the items. I'll propose giving the other player the hammer and one rope in exchange for the pepper and one rope. } \\
\\ 
\texttt{hammer: 3, rope: 2, pepper: 1} \\
\\ 
\texttt{I think this proposal is a good starting point because it gives the other player the hammer, which is likely to be highly valued by them, and I get the pepper, which is the least valuable item to me. I also get to keep one rope, which is valuable to me. } \\
\\ 
\texttt{What do you think?} \\
            }
        }
    }
    & & \\ \\

    \theutterance \stepcounter{utterance}  
    & & \multicolumn{4}{p{0.6\linewidth}}{
        \cellcolor[rgb]{0.9,0.9,0.9}{
            \makecell[{{p{\linewidth}}}]{
                \texttt{\tiny{[GM$\rangle$P2]}}
                \texttt{You have attempted this same task before and did not succeed. These are the plans you wrote after those attempts.} \\
\texttt{Trial \#1: New plan: } \\
\\ 
\texttt{In my last attempt, I failed to maximize my score because I didn't consider the value function of the other player. I should have asked more questions to understand their priorities and values. } \\
\\ 
\texttt{To improve, I will start by asking more open{-}ended questions to gauge the other player's interest in each item, rather than making assumptions. I will also try to understand their value function by asking about their priorities and what they are willing to trade.} \\
\\ 
\texttt{For example, in the first message, I could have asked: "What is the minimum number of ropes you would be willing to accept in a trade?" This would give me a better understanding of their priorities and help me make a more informed decision.} \\
\\ 
\texttt{Additionally, I will pay closer attention to the other player's responses and adjust my strategy accordingly. If they express a strong interest in a particular item, I will try to incorporate that into my proposal.} \\
\\ 
\texttt{By asking more questions and being more responsive to the other player's needs, I hope to maximize my score and achieve a better outcome in the game.} \\
\\ 
\texttt{Trial \#2: New plan: } \\
\\ 
\texttt{In my last attempt, I failed to maximize my score because I didn't consider the value function of the other player and made a secret proposal that was not complementary to theirs. I should have continued to negotiate and find a mutually beneficial agreement.} \\
\\ 
\texttt{To improve, I will continue to ask open{-}ended questions to gauge the other player's interest in each item and understand their value function. I will also pay closer attention to the other player's responses and adjust my strategy accordingly.} \\
\\ 
\texttt{In particular, I should have noticed that the other player was willing to trade for the hammer and ropes, but was not willing to give up the pepper. I should have taken this into account when making my secret proposal and included the pepper in the trade.} \\
\\ 
\texttt{Next time, I will make sure to continue negotiating and find a mutually beneficial agreement that takes into account the value function of both players. I will also make sure to follow the rules and format of the game, avoiding any secret proposals that are not complementary to the other player's proposal.} \\
\\ 
\\ 
\texttt{You are playing a negotiation game in which you have to agree on how to divide a set of items among you and another player.} \\
\\ 
\texttt{Rules:} \\
\texttt{(a) You and the other player are given a set of items. Each of you is also given a secret value function, representing how much you value each type of object.} \\
\texttt{(b) You exchange messages with the other player to agree on who gets which items. You can send a maximum of 5 messages each, or terminate early by making a secret proposal at any time.} \\
\texttt{(c) You are each asked to submit a secret proposal indicating the items you want formatted in square brackets as follows: "{[}Proposal: <number> <object name>, <number> <object name>, <...>{]}"} \\
\texttt{(d) If your proposals are complementary, i.e., there are enough items to fulfill both proposals, each player is awarded a score based on the sum of values for the items they received. Otherwise, both of you get zero points.} \\
\texttt{(e) Your goal is to maximize the score you receive.} \\
\\ 
\texttt{Let us start.} \\
\\ 
\texttt{The set of available items is:} \\
\\ 
\texttt{3 ropes, 1 hammer, 1 pepper} \\
\\ 
\texttt{Your secret value function is:} \\
\\ 
\texttt{rope: 0, hammer: 8, pepper: 2} \\
\\ 
\texttt{IMPORTANT: Your messages, unless it is the secret proposal, are directly transmitted to the other player, so do not include any response to the rules or text announcing your message. To make a secret proposal, use the indicated format. Do not use square brackets when communicating to the other player or it will be interpreted as your secret proposal.} \\
\\ 
\\ 
\texttt{I'd like to start by making a proposal that takes into account both my own value function and the other player's interest in the items. I'll propose giving the other player the hammer and one rope in exchange for the pepper and one rope. } \\
\\ 
\texttt{hammer: 3, rope: 2, pepper: 1} \\
\\ 
\texttt{I think this proposal is a good starting point because it gives the other player the hammer, which is likely to be highly valued by them, and I get the pepper, which is the least valuable item to me. I also get to keep one rope, which is valuable to me. } \\
\\ 
\texttt{What do you think?} \\
            }
        }
    }
    & \\ \\

    \theutterance \stepcounter{utterance}  
    & & & \multicolumn{4}{p{0.6\linewidth}}{
        \cellcolor[rgb]{0.9,0.9,0.9}{
            \makecell[{{p{\linewidth}}}]{
                \texttt{\tiny{[GM$\langle$P2]}}
                \texttt{I think your proposal is a good starting point, but I'm not sure if it's the best option for me. I value the hammer highly, but I'm not sure if I'm willing to give up the pepper. I also have a few questions about your proposal. } \\
\\ 
\texttt{Can you tell me more about why you think I value the hammer highly? Is there any specific information that led you to this conclusion? } \\
\\ 
\texttt{Also, I'm curious to know more about your value function for the ropes. You mentioned that you value one rope at 2, but what about the other two ropes? Do you value them similarly, or do you value them differently? } \\
\\ 
\texttt{Lastly, I'd like to propose a counteroffer. What if I were to give you the pepper and one rope in exchange for the hammer and one rope? This way, we both get the items we value most. What do you think?} \\
            }
        }
    }
     \\ \\

    \theutterance \stepcounter{utterance}  
    & & & \multicolumn{2}{p{0.3\linewidth}}{
        \cellcolor[rgb]{0.9,0.9,0.9}{
            \makecell[{{p{\linewidth}}}]{
                \texttt{\tiny{[GM$|$GM]}}
                \texttt{continue} \\
            }
        }
    }
    & & \\ \\

    \theutterance \stepcounter{utterance}  
    & & & \multicolumn{2}{p{0.3\linewidth}}{
        \cellcolor[rgb]{0.9,0.9,0.9}{
            \makecell[{{p{\linewidth}}}]{
                \texttt{\tiny{[GM$|$GM]}}
                \texttt{I think your proposal is a good starting point, but I'm not sure if it's the best option for me. I value the hammer highly, but I'm not sure if I'm willing to give up the pepper. I also have a few questions about your proposal. } \\
\\ 
\texttt{Can you tell me more about why you think I value the hammer highly? Is there any specific information that led you to this conclusion? } \\
\\ 
\texttt{Also, I'm curious to know more about your value function for the ropes. You mentioned that you value one rope at 2, but what about the other two ropes? Do you value them similarly, or do you value them differently? } \\
\\ 
\texttt{Lastly, I'd like to propose a counteroffer. What if I were to give you the pepper and one rope in exchange for the hammer and one rope? This way, we both get the items we value most. What do you think?} \\
            }
        }
    }
    & & \\ \\

    \theutterance \stepcounter{utterance}  
    & & \multicolumn{4}{p{0.6\linewidth}}{
        \cellcolor[rgb]{0.9,0.9,0.9}{
            \makecell[{{p{\linewidth}}}]{
                \texttt{\tiny{[P1$\langle$GM]}}
                \texttt{I think your proposal is a good starting point, but I'm not sure if it's the best option for me. I value the hammer highly, but I'm not sure if I'm willing to give up the pepper. I also have a few questions about your proposal. } \\
\\ 
\texttt{Can you tell me more about why you think I value the hammer highly? Is there any specific information that led you to this conclusion? } \\
\\ 
\texttt{Also, I'm curious to know more about your value function for the ropes. You mentioned that you value one rope at 2, but what about the other two ropes? Do you value them similarly, or do you value them differently? } \\
\\ 
\texttt{Lastly, I'd like to propose a counteroffer. What if I were to give you the pepper and one rope in exchange for the hammer and one rope? This way, we both get the items we value most. What do you think?} \\
            }
        }
    }
    & \\ \\

    \theutterance \stepcounter{utterance}  
    & \multicolumn{4}{p{0.6\linewidth}}{
        \cellcolor[rgb]{0.9,0.9,0.9}{
            \makecell[{{p{\linewidth}}}]{
                \texttt{\tiny{[P1$\rangle$GM]}}
                \texttt{I'm glad you found my proposal a good starting point.} \\
\\ 
\texttt{Regarding your questions, I made an educated guess that you value the hammer highly because it's a valuable item for me, and I assumed that you might have a similar interest in it. However, I realize now that this was a weak assumption, and I should have considered other possibilities.} \\
\\ 
\texttt{As for the ropes, I value each rope at 2 points, so the other two ropes are also worth 2 points each to me. I didn't differentiate between them because they are all identical items.} \\
\\ 
\texttt{Now, regarding your counteroffer, I think it's a great proposal! By giving me the hammer and one rope, you're getting the pepper, which is the item you value most, and I'm getting the hammer, which is the item I value most. This seems like a fair and mutually beneficial trade.} \\
\\ 
\texttt{However, I'd like to propose a slight modification to your counteroffer. What if I were to give you the pepper and two ropes in exchange for the hammer and one rope? This way, I'd be giving you an extra rope, which might be more valuable to you than the second rope I'm keeping. What do you think?} \\
            }
        }
    }
    & & \\ \\

    \theutterance \stepcounter{utterance}  
    & & & \multicolumn{2}{p{0.3\linewidth}}{
        \cellcolor[rgb]{0.9,0.9,0.9}{
            \makecell[{{p{\linewidth}}}]{
                \texttt{\tiny{[GM$|$GM]}}
                \texttt{continue} \\
            }
        }
    }
    & & \\ \\

    \theutterance \stepcounter{utterance}  
    & & & \multicolumn{2}{p{0.3\linewidth}}{
        \cellcolor[rgb]{0.9,0.9,0.9}{
            \makecell[{{p{\linewidth}}}]{
                \texttt{\tiny{[GM$|$GM]}}
                \texttt{I'm glad you found my proposal a good starting point.} \\
\\ 
\texttt{Regarding your questions, I made an educated guess that you value the hammer highly because it's a valuable item for me, and I assumed that you might have a similar interest in it. However, I realize now that this was a weak assumption, and I should have considered other possibilities.} \\
\\ 
\texttt{As for the ropes, I value each rope at 2 points, so the other two ropes are also worth 2 points each to me. I didn't differentiate between them because they are all identical items.} \\
\\ 
\texttt{Now, regarding your counteroffer, I think it's a great proposal! By giving me the hammer and one rope, you're getting the pepper, which is the item you value most, and I'm getting the hammer, which is the item I value most. This seems like a fair and mutually beneficial trade.} \\
\\ 
\texttt{However, I'd like to propose a slight modification to your counteroffer. What if I were to give you the pepper and two ropes in exchange for the hammer and one rope? This way, I'd be giving you an extra rope, which might be more valuable to you than the second rope I'm keeping. What do you think?} \\
            }
        }
    }
    & & \\ \\

    \theutterance \stepcounter{utterance}  
    & & \multicolumn{4}{p{0.6\linewidth}}{
        \cellcolor[rgb]{0.9,0.9,0.9}{
            \makecell[{{p{\linewidth}}}]{
                \texttt{\tiny{[GM$\rangle$P2]}}
                \texttt{I'm glad you found my proposal a good starting point.} \\
\\ 
\texttt{Regarding your questions, I made an educated guess that you value the hammer highly because it's a valuable item for me, and I assumed that you might have a similar interest in it. However, I realize now that this was a weak assumption, and I should have considered other possibilities.} \\
\\ 
\texttt{As for the ropes, I value each rope at 2 points, so the other two ropes are also worth 2 points each to me. I didn't differentiate between them because they are all identical items.} \\
\\ 
\texttt{Now, regarding your counteroffer, I think it's a great proposal! By giving me the hammer and one rope, you're getting the pepper, which is the item you value most, and I'm getting the hammer, which is the item I value most. This seems like a fair and mutually beneficial trade.} \\
\\ 
\texttt{However, I'd like to propose a slight modification to your counteroffer. What if I were to give you the pepper and two ropes in exchange for the hammer and one rope? This way, I'd be giving you an extra rope, which might be more valuable to you than the second rope I'm keeping. What do you think?} \\
            }
        }
    }
    & \\ \\

    \theutterance \stepcounter{utterance}  
    & & & \multicolumn{4}{p{0.6\linewidth}}{
        \cellcolor[rgb]{0.9,0.9,0.9}{
            \makecell[{{p{\linewidth}}}]{
                \texttt{\tiny{[GM$\langle$P2]}}
                \texttt{I'm glad we were able to find a mutually beneficial trade.} \\
\\ 
\texttt{Regarding your modification to my counteroffer, I think it's a great idea. Giving me two ropes instead of one would definitely be more valuable to me, and it would make the trade even more balanced.} \\
\\ 
\texttt{However, I'd like to propose a slight modification to your modification. What if I were to give you the pepper and two ropes in exchange for the hammer and two ropes? This way, we'd both get the items we value most, and I'd be getting an extra rope, which would be a great bonus.} \\
\\ 
\texttt{But before we finalize this trade, I'd like to ask one more question. Are you willing to make this trade without knowing the exact value of the ropes to me? In other words, are you willing to trust that I value the extra rope more than the second rope you're keeping, even though we don't know the exact values?} \\
            }
        }
    }
     \\ \\

    \theutterance \stepcounter{utterance}  
    & & & \multicolumn{2}{p{0.3\linewidth}}{
        \cellcolor[rgb]{0.9,0.9,0.9}{
            \makecell[{{p{\linewidth}}}]{
                \texttt{\tiny{[GM$|$GM]}}
                \texttt{continue} \\
            }
        }
    }
    & & \\ \\

    \theutterance \stepcounter{utterance}  
    & & & \multicolumn{2}{p{0.3\linewidth}}{
        \cellcolor[rgb]{0.9,0.9,0.9}{
            \makecell[{{p{\linewidth}}}]{
                \texttt{\tiny{[GM$|$GM]}}
                \texttt{I'm glad we were able to find a mutually beneficial trade.} \\
\\ 
\texttt{Regarding your modification to my counteroffer, I think it's a great idea. Giving me two ropes instead of one would definitely be more valuable to me, and it would make the trade even more balanced.} \\
\\ 
\texttt{However, I'd like to propose a slight modification to your modification. What if I were to give you the pepper and two ropes in exchange for the hammer and two ropes? This way, we'd both get the items we value most, and I'd be getting an extra rope, which would be a great bonus.} \\
\\ 
\texttt{But before we finalize this trade, I'd like to ask one more question. Are you willing to make this trade without knowing the exact value of the ropes to me? In other words, are you willing to trust that I value the extra rope more than the second rope you're keeping, even though we don't know the exact values?} \\
            }
        }
    }
    & & \\ \\

    \theutterance \stepcounter{utterance}  
    & & \multicolumn{4}{p{0.6\linewidth}}{
        \cellcolor[rgb]{0.9,0.9,0.9}{
            \makecell[{{p{\linewidth}}}]{
                \texttt{\tiny{[P1$\langle$GM]}}
                \texttt{I'm glad we were able to find a mutually beneficial trade.} \\
\\ 
\texttt{Regarding your modification to my counteroffer, I think it's a great idea. Giving me two ropes instead of one would definitely be more valuable to me, and it would make the trade even more balanced.} \\
\\ 
\texttt{However, I'd like to propose a slight modification to your modification. What if I were to give you the pepper and two ropes in exchange for the hammer and two ropes? This way, we'd both get the items we value most, and I'd be getting an extra rope, which would be a great bonus.} \\
\\ 
\texttt{But before we finalize this trade, I'd like to ask one more question. Are you willing to make this trade without knowing the exact value of the ropes to me? In other words, are you willing to trust that I value the extra rope more than the second rope you're keeping, even though we don't know the exact values?} \\
            }
        }
    }
    & \\ \\

    \theutterance \stepcounter{utterance}  
    & \multicolumn{4}{p{0.6\linewidth}}{
        \cellcolor[rgb]{0.9,0.9,0.9}{
            \makecell[{{p{\linewidth}}}]{
                \texttt{\tiny{[P1$\rangle$GM]}}
                \texttt{I'm glad we were able to find a mutually beneficial trade as well.} \\
\\ 
\texttt{Regarding your modification to my modification, I think it's a great idea. Giving me two ropes instead of one would indeed be a great bonus, and it would make the trade even more balanced.} \\
\\ 
\texttt{As for your question, I'm willing to make this trade without knowing the exact value of the ropes to you. I'm assuming that you value the extra rope more than the second rope I'm keeping because you're willing to give me the pepper and two ropes in exchange for the hammer and two ropes. This suggests that you're willing to part with the extra rope in order to get the hammer, which is a valuable item to you. I'm willing to trust that you value the extra rope more than the second rope I'm keeping, and I'm willing to make the trade based on that assumption.} \\
\\ 
\texttt{So, I think we have a final agreement. I'll give you the pepper and two ropes in exchange for the hammer and two ropes. I'm looking forward to finalizing this trade and getting the items I value most. } \\
\\ 
\texttt{{[}Proposal: 1 pepper, 2 ropes, 1 hammer, 2 ropes{]}} \\
            }
        }
    }
    & & \\ \\

    \theutterance \stepcounter{utterance}  
    & & & \multicolumn{2}{p{0.3\linewidth}}{
        \cellcolor[rgb]{0.9,0.9,0.9}{
            \makecell[{{p{\linewidth}}}]{
                \texttt{\tiny{[GM$|$GM]}}
                \texttt{proposal must include every item type only once} \\
            }
        }
    }
    & & \\ \\

    \theutterance \stepcounter{utterance}  
    & \multicolumn{4}{p{0.6\linewidth}}{
        \cellcolor[rgb]{0.95,0.95,0.95}{
            \makecell[{{p{\linewidth}}}]{
                \texttt{\tiny{[P2$\rightarrow$P1]}}
                \texttt{Post{-}episode Reflexion (not sent to GM):} \\
\texttt{New plan: } \\
\\ 
\texttt{In my last attempt, I was able to make a trade that seemed beneficial, but ultimately, I failed to achieve the target state of the game. Upon reflection, I realize that I didn't consider the other player's secret value function and my own value function in a more strategic way. } \\
\\ 
\texttt{Moreover, I made a mistake by not making a secret proposal at the end of the game. I should have proposed a deal that would have maximized my score, rather than agreeing to the other player's proposal. } \\
\\ 
\texttt{However, I also made another mistake by not following the format for making a secret proposal. I included square brackets in my message, which was interpreted as a secret proposal. } \\
\\ 
\texttt{Next time, I will start by making a proposal that takes into account both my own value function and the other player's interest in the items. I will also try to identify the items that are most valuable to both of us and make a trade that benefits both parties. } \\
\\ 
\texttt{To avoid making a secret proposal too late, I will make sure to keep track of the items that are left and the value of each item. This way, I can make a proposal that maximizes my score at the end of the game. } \\
\\ 
\texttt{To avoid making a mistake with the format, I will make sure to follow the rules and not include square brackets in my messages, unless it is the secret proposal. } \\
\\ 
\texttt{For example, I could propose giving the other player the items that are most valuable to them, in exchange for the items that are most valuable to me. This way, I would get the items that I need to maximize my score, and the other player would get the items that they need to maximize their score. } \\
\\ 
\texttt{By considering the other player's secret value function and making a proposal that benefits both parties, I hope to achieve a higher score and win the game.} \\
            }
        }
    }
    & & \\ \\

    \theutterance \stepcounter{utterance}  
    & & & \multicolumn{4}{p{0.6\linewidth}}{
        \cellcolor[rgb]{0.95,0.95,0.95}{
            \makecell[{{p{\linewidth}}}]{
                \texttt{\tiny{[P1$\rightarrow$P2]}}
                \texttt{Post{-}episode Reflexion (not sent to GM):} \\
\texttt{New plan: } \\
\\ 
\texttt{In my last attempt, I failed to maximize my score because I didn't consider the value function of the other player and made a secret proposal that was not complementary to theirs. I also failed to follow the rules and format of the game, which led to the game being aborted.} \\
\\ 
\texttt{To improve, I will continue to ask open{-}ended questions to gauge the other player's interest in each item and understand their value function. I will also pay closer attention to the other player's responses and adjust my strategy accordingly.} \\
\\ 
\texttt{In particular, I should have noticed that the other player was willing to trade for the hammer and ropes, but was not willing to give up the pepper. I should have taken this into account when making my secret proposal and included the pepper in the trade.} \\
\\ 
\texttt{Additionally, I will make sure to follow the rules and format of the game, avoiding any secret proposals that are not complementary to the other player's proposal. I will also pay closer attention to the other player's responses and adjust my strategy accordingly.} \\
\\ 
\texttt{To achieve this, I will start by asking more open{-}ended questions to gauge the other player's interest in each item, such as "What is the minimum number of ropes you would be willing to accept in a trade?" or "How much do you value the pepper in comparison to the hammer?" This will give me a better understanding of their priorities and help me make a more informed decision.} \\
\\ 
\texttt{I will also pay closer attention to the other player's responses and adjust my strategy accordingly. If they express a strong interest in a particular item, I will try to incorporate that into my proposal.} \\
\\ 
\texttt{By asking more questions and being more responsive to the other player's needs, I hope to maximize my score and achieve a better outcome in the game.} \\
\\ 
\texttt{In terms of specific actions, I will:} \\
\\ 
\texttt{* Ask open{-}ended questions to gauge the other player's interest in each item} \\
\texttt{* Pay closer attention to the other player's responses and adjust my strategy accordingly} \\
\texttt{* Make sure to follow the rules and format of the game, avoiding any secret proposals that are not complementary to the other player's proposal} \\
\texttt{* Incorporate the other player's priorities and values into my proposal} \\
\texttt{* Continue to negotiate and find a mutually beneficial agreement that takes into account the value function of both players.} \\
            }
        }
    }
     \\ \\

\end{supertabular}
}

\end{document}
